\section{第10章の解答}\label{解答:第10章}

\paragraph{10.1}
分母から$n^2$をくくり,分子の$k$を$n\cdot\frac{k}{n}$と書くと
\[
 \frac{k}{n^2+k^2}
 =\frac1n\cfrac{\cfrac{k}{n}}{1+\left(\cfrac{k}{n}\right)^2}.
\]
したがって,連続関数$f(x)=\frac{x}{1+x^2}$の$[0,1]$上の区分求積であり,求める極限は
\[
 \int_0^1\frac{x}{1+x^2}\dd x
 =\left[\frac12\log(1+x^2)\right]_0^1
 =\frac12\log2.
\]

\paragraph{10.2}
$\sqrt{n+k}=\sqrt n\sqrt{1+\frac{k}{n}}$なので,もとの和は
\[
 \frac1n\sum_{k=1}^{n}\sqrt{1+\frac{k}{n}}
\]
となる.\,$f(x)=\sqrt{1+x}$は$[0,1]$で連続だから,極限は
\[
 \int_0^1\sqrt{1+x}\dd x
 =\left[\frac23(1+x)^{\frac32}\right]_0^1
 =\frac23(2\sqrt2-1).
\]
分子の根号から出る$\sqrt n$と,前にある$n^{\frac32}$を合わせて,区分求積の幅$\frac1n$が現れる.

\paragraph{10.3}
$f(x)=\frac1{x^2}$は$x>0$で減少するので,\,$k\geqq2$に対して
\[
 \int_k^{k+1}\frac1{x^2}\dd x
 \leqq\frac1{k^2}
 \leqq\int_{k-1}^{k}\frac1{x^2}\dd x.
\]
$k=n+1,\ldots,N$について足せば
\[
 \int_{n+1}^{N+1}\frac1{x^2}\dd x
 \leqq\sum_{k=n+1}^{N}\frac1{k^2}
 \leqq\int_n^N\frac1{x^2}\dd x.
\]
原始関数は$-\frac1x$だから,問題の不等式を得る.

部分和$S_N=\sum_{k=1}^{N}\frac1{k^2}$は単調増加する.一方,\,$n=1$の右側の不等式を使うと
\[
 S_N\leqq1+1-\frac1N<2\qquad(N\geqq2)
\]
なので,上に有界である.したがって有限の極限$S$を持つ.これを無限級数の和とし,\,$R_n=S-S_n$とおく.最初の不等式で$n$を固定して$N\to\infty$とすると
\[
 \frac1{n+1}\leqq R_n\leqq\frac1n
\]
を得る.収束の証明では,和の正確な値を求める必要はなかった.

\paragraph{10.4}
コーシー・シュワルツの不等式を$f(x)$と$x$に使うと
\[
 1=\left(\int_0^1xf(x)\dd x\right)^2
 \leqq\left(\int_0^1f(x)^2\dd x\right)
       \left(\int_0^1x^2\dd x\right)
 =\frac13\int_0^1f(x)^2\dd x.
\]
よって積分は$3$以上である.等号条件は$f(x)=cx$であり,制約に代入すると
\[
 1=\int_0^1cx^2\dd x=\frac c3
\]
から$c=3$となる.実際$f(x)=3x$なら制約を満たし,\,$\int_0^1(3x)^2\dd x=3$である.したがって最小値は$3$,それを与える関数は$f(x)=3x$である.

\paragraph{10.5}
まず$K_0=\int_1^{\e}1\dd x=\e-1$である.\,$n\geqq1$では$(\log x)^n$を微分する部分積分によって
\begin{align*}
 K_n
 &=\left[x(\log x)^n\right]_1^{\e}
   -n\int_1^{\e}x(\log x)^{n-1}\frac1x\dd x\\
 &=\e-nK_{n-1}.
\end{align*}
したがって
\[
 K_1=1,\qquad K_2=\e-2,\qquad K_3=\e-3(\e-2)=6-2\e.
\]
$n=0$を漸化式に形式的に代入すると$K_0=\e$という誤りが生じる.\,$x=1$の端点の項が$0$になるという議論は$n\geqq1$に対して行ったものである.

\paragraph{10.6}
本文の$I_n=\int_0^{\frac{\pi}{2}}\sin^n x\dd x$について
\[
 \frac{2m+1}{2m+2}
 \leqq\frac{I_{2m+1}}{I_{2m}}\leqq1,
 \qquad
 P_m=\frac{\pi}{2}\frac{I_{2m+1}}{I_{2m}}
\]
である.したがって
\[
 \frac{\pi}{2}\frac{2m+1}{2m+2}
 \leqq P_m\leqq\frac{\pi}{2}.
\]
$\frac{\pi}{2}$から各辺を引くと不等号の向きが逆になり
\[
 0\leqq\frac{\pi}{2}-P_m
 \leqq\frac{\pi}{2}\left(1-\frac{2m+1}{2m+2}\right)
 =\frac{\pi}{4m+4}.
\]
次に$u=\frac{\pi}{2}-x$とおけば
\[
 \int_0^{\frac{\pi}{2}}\cos^{2m}x\dd x=I_{2m}.
\]
本文で$I_{2m}I_{2m+1}=\frac{\pi}{2(2m+1)}$および$\frac{I_{2m+1}}{I_{2m}}\to1$を示したので,
\[
 mI_{2m}^2
 =\frac{m\pi}{2(2m+1)}\frac{I_{2m}}{I_{2m+1}}
 \longrightarrow\frac{\pi}{4}.
\]
積分値が正であることを使って平方根を取り,求める極限は$\frac{\sqrt\pi}{2}$となる.

\paragraph{10.7}
整数の場合のベータ積分を使うと
\begin{align*}
 \int_0^1\binom{m+n}{m}x^m(1-x)^n\dd x
 &=\frac{(m+n)!}{m!n!}\frac{m!n!}{(m+n+1)!}\\
 &=\frac1{m+n+1}.
\end{align*}
また$m=2,n=4$の場合は
\[
 \int_0^1x^2(1-x)^4\dd x
 =\frac{2!4!}{7!}=\frac1{105}.
\]

\paragraph{10.8}
$0<x\leqq\frac12$では$1-x\geqq\frac12$だから
\[
 0\leqq\sqrt{\frac{x}{1-x}}\leqq\sqrt2\sqrt x.
\]
右辺は$0$まで積分できる.\,$\frac12\leqq x<1$では
\[
 0\leqq\sqrt{\frac{x}{1-x}}\leqq\frac1{\sqrt{1-x}},
\]
右辺の積分は$[-2\sqrt{1-x}]_{\frac12}^{1}=\sqrt2$で有限である.各端点に近づけた積分値は単調増加し,これらの有限な値で上から抑えられるので,両端で収束する.

被積分関数は$x^{\frac32-1}(1-x)^{\frac12-1}$なので,積分は$B\left(\frac32,\frac12\right)$である.漸化式より
\[
 B\left(\frac32,\frac12\right)
 =\frac12 B\left(\frac12,\frac12\right)=\frac{\pi}{2}.
\]
置換$x=\sin^2t$でも,開区間では平方根が$\tan t$となり,
\[
 \int_0^1\sqrt{\frac{x}{1-x}}\dd x
 =2\int_0^{\frac{\pi}{2}}\sin^2t\dd t
 =2\cdot\frac{\pi}{4}=\frac{\pi}{2}
\]
と一致する.置換は端点を避けた区間で行い,その後で極限を取る.
